% Owned by the skeleton thread.
\section{Reproducing this report}
\label{sec:reproducibility}

This document was built from commit \commit{\buildcommit} of
\url{https://github.com/vehagn/vibe-rations}\ifbuilddirty, \textcolor{vrred}{with
uncommitted changes in the working tree}\fi. Every number, table and figure that
comes from the solver is the product of an \emph{experiment}: a short manifest under
\file{report/experiments} naming the commands to run. This appendix says how to rerun them,
how the report knows whether its numbers are those of the code it was built from, and lists
every experiment with its commands.

\subsection{Building the report}
\label{sec:reproducibility:build}

With the repository checked out and its development container available
(\cref{sec:usage:start}),
\begin{lstlisting}[style=shell]
just in report                # the PDF from the committed data: report/build/vibe-rations-report.pdf
just in report-status         # every experiment and whether its data are current
just in report-run            # rerun the quick experiments and record their provenance
just in report-run --stale --tier quick --tier study   # rerun what has gone stale
just in report-run verification/stokes                 # rerun one experiment by name
\end{lstlisting}
\code{just report} needs no simulation: the experiments' results are committed, reduced to
small text files, under \file{report/data}, and the figures are drawn from them by LaTeX
itself (pgfplots) each time the PDF is built.

\subsection{How a number gets here}
\label{sec:reproducibility:how}

\begin{enumerate}
  \item \textbf{A manifest names the commands.} An experiment is a file
    \code{report/experiments/<section>/<name>.toml} with a title, a tier and a list of shell
    commands, run from the repository's root. The commands are the ones a person would type:
    the solver's own subcommands (\cref{sec:usage}), the tests that compare it with closed
    forms, and short scripts that reduce what those wrote. A script or a directory of inputs
    that belongs to the experiment sits beside its manifest under the same name.
  \item \textbf{The generator runs them and records what it did.}
    \code{report/report.py run <section>/<name>} builds the solver, empties the experiment's
    data directory \code{report/data/<section>/<name>}, runs the commands in order, and then
    writes \code{provenance.json} there: the commit checked out, any uncommitted change in
    what the experiment depends on, the content hash of each thing it depends on, every
    command with the time it took, the machine (its processor, threads and memory; no host
    name), the compiler, and the size and SHA-256 of every file produced. A command that
    fails leaves no provenance, so half a result can't pass for a whole one.
  \item \textbf{Only reduced results are committed.} The data directory holds small text
    files: curves thinned to a few hundred points as CSV, single numbers as
    \code{values.json}, a small PNG where a field is shown. Whatever is bulky (whole
    wavefields, every trace) stays in the git-ignored \code{report/build/raw}. Because the
    committed files are text with six significant digits, a rerun after a change to the
    solver shows up in \code{git diff} as the numbers that moved.
  \item \textbf{LaTeX draws from the data.} Figures are pgfplots reading the CSV files, tables
    are typeset from them, and a number in the text is \code{\val{<section>/<name>/<key>}},
    looked up in \code{values.json}. Nothing between the committed data and the PDF is drawn
    or typed by hand, and building the PDF runs no simulation.
  \item \textbf{The PDF says where each result is from.} A caption ends with the
    experiment's name and the commit its data were made at, linked to its entry below.
\end{enumerate}

\paragraph{Current, stale, dirty, missing.} The results are committed, so they can fall
behind the code. The generator therefore compares, for each experiment, the content hashes
recorded in its provenance with those of the commit the report is built from: the solver's
sources (\code{src}), \code{Cargo.toml}, \code{Cargo.lock}, \code{rust-toolchain.toml}, the
experiment's manifest, scripts and inputs, and whatever else the manifest lists (a test
file, an example's configuration). If all agree the experiment is \emph{current}: its
numbers are those of this commit's code, even though they were computed at an earlier commit.
If any differs it is \emph{stale}; if it was run with uncommitted changes in those files it
is \emph{dirty}; with no data at all it is \emph{missing}. Anything but current is marked in
red wherever the experiment is used and in \cref{tab:reproducibility:experiments}, and
\code{report/report.py check} fails on it. The working order is therefore: commit the code,
run the experiment, commit its data.

\paragraph{Tiers.} Experiments are sorted by cost, so that keeping the report current
doesn't mean rerunning everything after every change: \code{quick} ones take up to a couple
of minutes each, \code{study} ones up to an hour, and \code{heavy} ones longer, or need a
large machine or downloaded data (their manifests say which). \code{just report-run} alone
runs the quick ones.

\paragraph{What reproduces exactly.} The solver is deterministic: the same build on the
same machine gives the same traces whatever the number of threads and whichever compiled
copy of the update loops runs. The experiment \code{reproducibility/determinism} checks that
each time it is run, by making the tutorial's run of \cref{sec:usage:explosion} three times
(with every thread, with one, and with the plain loops in place of the
\texttt{\val{reproducibility/determinism/kernels_default}} ones) and comparing the SHA-256 of
the trace files: the \val{reproducibility/determinism/files} files were
\textbf{\val{reproducibility/determinism/identical}} (\cref{tab:reproducibility:determinism}).
A rerun of an experiment therefore reproduces its committed files byte for byte, apart from
\code{provenance.json}'s dates and durations and any output that prints a speed.

\begin{table}[htb]
  \centering
  \caption{The first 16 hexadecimal digits of the SHA-256 of each trace file of the same
    run made three ways. \source{reproducibility/determinism}}
  \label{tab:reproducibility:determinism}
  \pgfplotstabletypeset[
    string type, font=\small\ttfamily,
    columns/file/.style={column name={\normalfont File}, column type=l},
    columns/threads-all/.style={column name={\normalfont All threads}},
    columns/threads-one/.style={column name={\normalfont One thread}},
    columns/plain-loops/.style={column name={\normalfont Plain loops}},
    columns/identical/.style={column name={\normalfont Identical}},
  ]{\data/reproducibility/determinism/hashes.csv}
\end{table}

Across machines the last digits can differ. The repository's notes (\file{AGENTS.md}) record
one such case: between Linux and macOS the time step differed in its last bit, most likely
through the maths library used when the stencil's coefficients are designed, and traces then
agreed to $6 \times 10^{-12}$ of their peak rather than bit for bit. The six significant
digits of the committed files hide most of that, but a rerun elsewhere should be expected to
match to about that many digits, not to produce identical files.

\subsection{The experiments}
\label{sec:reproducibility:experiments}

\Cref{tab:reproducibility:experiments} lists every experiment, the commit its data were made
at, and whether the code it depends on has changed since. The commands follow.

\begin{longtable}{@{}llllll@{}}
  \caption{The report's experiments.\label{tab:reproducibility:experiments}}\\
  \toprule
  Experiment & Tier & Commit & Made & Took & Status\\
  \midrule
  \endfirsthead
  \toprule
  Experiment & Tier & Commit & Made & Took & Status\\
  \midrule
  \endhead
  \bottomrule
  \endlastfoot
  \InputIfFileExists{\gen/table.tex}{}{\multicolumn{6}{c}{\missing{run report.py render}}\\}
\end{longtable}

\InputIfFileExists{\gen/appendix.tex}{}{}
